\begin{table}[htbp]
\centering
\caption{Espanha: indicadores fiscais por subsetor das administrações públicas, 2016--2025}
\label{tab:espanha}
\scriptsize
\begin{tabular}{lrrrrrrrr}
\toprule
 & \makecell{Receita\\\% PIB\\(média)} & \makecell{Despesa\\\% PIB\\(média)} & \makecell{Saldo\\\% PIB\\2016} & \makecell{Saldo\\\% PIB\\2025} & \makecell{Dívida\\\% PIB\\2016} & \makecell{Dívida\\\% PIB\\2025} & \makecell{Transf.\\recebidas/\\receita (\%)} & \makecell{Impostos\\próprios/\\receita (\%)}\\
\midrule
Governo geral & 40,7 & 45,0 & -4,2 & -2,4 & 102,0 & 100,5 & 4,9 & 57,0\\
Administração central & 18,4 & 21,5 & -2,5 & -2,0 & 89,9 & 92,5 & 11,0 & 78,5\\
Comunidades Autônomas & 16,0 & 16,5 & -0,8 & -0,4 & 24,7 & 20,2 & 58,9 & 35,4\\
Entidades locais & 6,5 & 6,1 & 0,6 & 0,3 & 2,9 & 1,2 & 44,3 & 47,4\\
Seguridade social & 14,7 & 15,7 & -1,5 & -0,3 & 1,5 & 8,1 & 15,8 & 0,0\\
\bottomrule
\end{tabular}
\par\smallskip
\begin{minipage}{0.97\textwidth}\scriptsize
Fonte: elaboração própria com dados do Eurostat (\texttt{gov\_10a\_main}, \texttt{gov\_10dd\_edpt1}, \texttt{nama\_10\_gdp}). Subsetores não consolidados entre si; transferências recebidas = D.7 + D.9 (inclui transferências intragovernamentais D.73 e D.92); impostos próprios = D.2 + D.5 + D.91, segundo o SEC 2010, que registra como receita regional a parcela cedida do IRPF e do IVA. Saldo = capacidade (+) ou necessidade ($-$) de financiamento (B.9).
\end{minipage}
\end{table}
